\section{Closest Work}\label{sec:related}
\subsection{The method under test}
scMINER \cite{pan2025scminer} introduces MICA and reports a mean adjusted Rand index
(ARI) of 0.84 across ten labelled datasets, against 0.73 for Seurat and 0.68 for Scanpy,
the two standard analysis toolkits. Every setting in that benchmark is fixed with the
labels, and nine of the ten datasets sit below the 5,000 cell threshold, so almost all of
the headline scores come from MICA given the true number of clusters. In the authors' own
benchmark files MICA trails Scanpy on the three larger datasets, which were used for
runtime rather than accuracy. Asked during peer review how a user should choose the
resolution without labels, the authors pointed to silhouette scores and marker genes but
tested neither, and MICA's shipped defaults appear only in its source code. We found no
independent benchmark that includes MICA.

\subsection{Choosing a clustering without labels}
chooseR \cite{pattersoncross2021chooser} is the rule we adopt. It reruns a pipeline on 100
random 80 percent subsamples per candidate setting, scores each setting by how reliably
cells co-cluster, and keeps the finest setting that stays above a threshold derived from
the settings themselves. Its authors note that preferring the finest acceptable setting
leans toward mild overclustering. It has not been applied to a mutual information method,
and the data it was demonstrated on, roughly 11,000 cells, is an order of magnitude larger
than nine of our ten datasets.

Two relatives return several answers instead of one: MultiK \cite{liu2021multik} pools
4,000 subsampled Seurat runs by cluster count, and Yu et al. \cite{yu2022benchmarking}
benchmarked 14 methods that estimate the number of cell types, finding that most
underestimate it and that a good ARI did not imply a correct count. That last point is why
we report the chosen count beside the score. A second family asks whether a split is real
rather than merely reproducible: sc-SHC \cite{grabski2023significance} tests each split
against a one-population null, returning one cluster on a single cell line where Seurat's
defaults returned nine, and CHOIR \cite{sant2025choir} overclusters then merges clusters a
random forest cannot separate. None of these covers a mutual information method, and the
distinction they draw matters here, because a partition can be perfectly reproducible
without being real.

\subsection{Tuning on the ground truth}
Ullmann et al. \cite{ullmann2023overoptimistic} tuned a published clustering algorithm
against known labels to make it look good, then showed how much of that advantage
disappeared when the choices moved. They argue the practice is common and overstates what
a user without labels obtains, and recommend tuning on separate datasets. Their study used
toy and UCI data and one algorithm, with no single-cell data and no label-free rule tested
as the remedy. A 2025 preprint \cite{szmigiel2025benchmarking} swept Leiden and Louvain
over a large grid on 16 single-cell datasets and found that on six the Scanpy default did
no better than a typical random setting, which shows the default is weak without studying
label-free selection either.

These studies establish the general concern: tuning on the labels flatters a method, and
shipped defaults are weak stand-ins for a choice. None measures it for MICA, whose published
scores were all obtained with label-informed settings. Our contribution is that measurement,
and with it something the general concern does not predict: the effect depends on which
cluster counts each method's grid can offer, so holding methods to a common cap on the number
of clusters is not enough to compare them.

\subsection{Scaling the neighbour graph} MICA's graph mode finds neighbours with the HNSW
graph index \cite{malkov2020efficient}. Proposing candidates cheaply and rescoring them
exactly is the standard alternative in nearest neighbour search, as in the refinement step
of Faiss \cite{douze2024faiss}, and like any shortlist it finds only the true neighbours that
land in it. Section~\ref{sec:graph} compares the two on cost and on the clustering each
yields. Song et al. \cite{song2012comparison} found that binned mutual information tracks
correlation closely for gene pairs, which is what suggested correlation as the cheap stage.
